\section{Introduction}
\label{sec:intro}

Two quantified AI scenarios now shape the UK debate, and this paper builds its two scenarios from them. The Office for Budget Responsibility's March~2026 \emph{Economic and fiscal outlook} includes a technological-displacement scenario in which AI substitutes for labour, deepens capital and lowers the labour share \citep[Box~2.2]{obr2026efo}. \citet{korinek2026} quantify a ``substantial'' US scenario to 2030 with a labour-share fall, a productivity gain and a capital-stock response. Both specify the factor-share and technology shocks. Neither works out how a UK economy with a calibrated tax system, an ageing population and an open capital account would respond to them in general equilibrium, or whose income the response raises.

This paper does that in OG-UK, the UK calibration of the OG-Core overlapping-generations model \citep{debacker2024ogcore}. The model has UN demographic projections, eighty annual cohorts with endogenous saving and labour supply, and effective tax functions estimated from PolicyEngine UK microdata. AI enters production
\begin{equation}
\label{eq:production}
Y_t = Z_t\, K_t^{\gamma_t} \bigl(e^{g_y t} L_t\bigr)^{1-\gamma_t}
\end{equation}
through the capital exponent $\gamma_t$, which carries automation, and the technology term $Z_t$. Both scenarios raise $\gamma$ linearly from \genGammaBase{} to \genGammaShocked{} over \genStartYear--\genEndYear{}, a \genLabourShareFallPp{} percentage-point fall in the labour share. \ShareShiftTfp{} also delivers the \genTfpGainPct{} per cent productivity gain of \citet{korinek2026}. \ShareShift{} has no productivity gain; Section~\ref{sec:calibration} gives its exact definition. The scenarios are imposed, not estimated. The question is what they imply once capital, labour and the public finances respond.

The paper makes three points.

\paragraph{The general-equilibrium factor response.} The shocks raise the return to capital, and capital rises in response. In \genEndYear{} the capital stock is \genShareShiftTfpCapitalGapThirtyCentralAbs{} per cent \genShareShiftTfpCapitalGapThirtyCentralDir{} the baseline in \ShareShiftTfp{} and \genShareShiftCapitalGapThirtyCentralAbs{} per cent \genShareShiftCapitalGapThirtyCentralDir{} it in \ShareShift{}. Aggregate labour input is \genShareShiftTfpLabourGapThirtyAbs{} per cent \genShareShiftTfpLabourGapThirtyDir{} the baseline in \ShareShiftTfp{} and \genShareShiftLabourGapThirtyAbs{} per cent \genShareShiftLabourGapThirtyDir{} it in \ShareShift{}.

\paragraph{National income, not only GDP.} How much of the extra capital is supplied from abroad depends on the capital closure. Under OG-UK's default ($\zetaK = \genZetaKCentral{}$, with $\zeta_D = \genZetaD{}$ and a \genWorldRatePct{} per cent world interest rate), about three-quarters of the extra 2030 capital is foreign-owned (\genShareShiftTfpForeignShareExtraKCentral{} and \genShareShiftForeignShareExtraKCentral{} per cent), so the return on it is paid abroad. I therefore lead with gross national income (GNI). In \genEndYear{}, GNI is \genShareShiftTfpGniGapThirtyCentralShort{} per cent against the baseline in \ShareShiftTfp{} and \genShareShiftGniGapThirtyCentralShort{} per cent in \ShareShift{}. The GDP gap is \genShareShiftTfpGdpGapThirtyCentralShort{} per cent in \ShareShiftTfp{}; in \ShareShift{} 2030 GDP is unchanged by construction. With private capital unable to flow in ($\zetaK = \genZetaKBound{}$) the GNI gaps are \genShareShiftTfpGniGapThirtyBoundShort{} and \genShareShiftGniGapThirtyBoundShort{} per cent. Both closures are reported together throughout (Table~\ref{tab:headline}). The value \genZetaKCentral{} is an OG-UK default, not a UK calibration; Section~\ref{sec:closure} sets out the after-the-fact case for using it and its weaknesses.

\paragraph{Wages.} At a unit elasticity of substitution the wage is $(1-\gamma_t)Y_t/L_t$, so it rises only if output per worker rises by more than the labour share falls. In \ShareShift{} it does not, and the real wage is \genShareShiftWageGapThirtyCentralAbs{} per cent \genShareShiftWageGapThirtyCentralDir{} the baseline in \genEndYear{} at $\zetaK = \genZetaKCentral{}$ and \genShareShiftWageGapThirtyBoundAbs{} per cent \genShareShiftWageGapThirtyBoundDir{} it at $\zetaK = \genZetaKBound{}$. In \ShareShiftTfp{} the two terms nearly cancel, and how much capital arrives decides the sign. The \genEndYear{} wage gap is \genShareShiftTfpWageGapThirtyBound{} per cent at $\zetaK = \genZetaKBound{}$ and \genShareShiftTfpWageGapThirtyCentral{} per cent at $\zetaK = \genZetaKCentral{}$, and at the central closure it turns positive in \genShareShiftTfpWageFirstPosYearCentral{} as capital keeps accumulating. The source scenario shows the same dependence: its average wage is \genKorinekTableFiveWageOne{} or \genKorinekTableFiveWageThree{} per cent depending on the elasticity of capital supply \citep[Table~5]{korinek2026}. The model does not determine whether wages rise or fall in that scenario.

Section~\ref{sec:incidence} uses the model's cohort structure to show who gains and who loses by birth cohort and lifetime-earnings group, a dimension that representative-agent and two-group AI models do not have. In \ShareShift{} at the central closure, cohorts near retirement in \genStartYear{} gain from the higher return on their assets (\genShareShiftCevAgeSixtyCentral{} per cent of consumption for those aged 60) and cohorts born from \genShareShiftCevFirstNegBirthCentral{} lose (\genShareShiftCevAgeTwentyCentral{} per cent for those aged 20). In \ShareShiftTfp{} every cohort gains. In both scenarios the bottom quarter of lifetime earners fares better than the top 1 per cent. These figures are consumption-equivalent changes in a utility that also counts leisure and bequests, and they depend on three features of the model that Section~\ref{sec:incidence} quantifies. Lump-sum transfers are the same amount per head and stay a fixed share of GDP while wages fall. The model's household wealth is spread more evenly across lifetime-earnings groups than earnings are. And in \ShareShift{} public borrowing is higher by \genShareShiftBorrowGapConsThirtyCentral{} per cent of consumption a year by \genEndYear{}, most of it received by households as lower taxes and higher interest on gilts, with no household paying for it inside the model. The incidence results are properties of the scenarios under those assumptions, not a finding that automation raises the living standards of lower earners.

All results are gaps against an OG-UK baseline that takes its productivity growth from the OBR's central forecast. Because the OBR's central forecast already assumes AI raises productivity growth \citep[paras~B.12 and B.14]{obr2025bp9}, the gaps are AI effects in addition to that assumption. Households have perfect foresight, so the scenarios move behaviour in \genStartYear{}, before technology has changed. I treat \genStartYear{} effects as anticipation, not as results. The model has no unemployment, no task reinstatement and no adoption friction. Section~\ref{sec:limitations} sets out these and the other boundaries.

The rest of the paper reviews the literature (Section~\ref{sec:literature}), sets out the model, the closure and the scenarios (Section~\ref{sec:methodology}), reports the aggregate results (Section~\ref{sec:results}) and the cohort incidence (Section~\ref{sec:incidence}), and discusses them (Section~\ref{sec:discussion}). Appendix~\ref{app:patch} documents the patch that lets OG-Core's capital exponent change over time, which the phased scenarios require.
